\documentclass[10pt]{article}
\usepackage[margin=0.7in]{geometry}
\usepackage[T1]{fontenc}
\usepackage[utf8]{inputenc}
\usepackage{lmodern}
\usepackage{microtype}
\usepackage{booktabs}
\usepackage{amsmath}
\usepackage{hyperref}
\usepackage{graphicx}

\hypersetup{colorlinks=true,linkcolor=blue,urlcolor=blue,citecolor=blue}
\setlength{\parskip}{0.3em}
\setlength{\parindent}{0pt}
\setlength{\textfloatsep}{6pt}
\setlength{\floatsep}{6pt}

\title{\vspace{-1.6em}\textbf{Synthetic RGB-D Pipeline for an Eye-in-Hand RealSense D455}\\
\large Modelled depth imperfections on aluminium profiles and brackets, and their effect on localization}
\author{[Your Name]}
\date{November 2026}

\begin{document}
\maketitle
\vspace{-1.2em}

\section{Task, approach, and scope}
A robot in a manufacturing cell picks and assembles aluminium profiles and brackets using an RGB-D camera. The task was to choose a commercial RGB-D camera, build a small pipeline that generates labelled synthetic RGB-D images of such a scene, model the camera's realistic imperfections instead of ideal depth, and explain which effects were included, which were left out, and why. I chose the \textbf{Intel RealSense D455}, an active IR stereo camera with a published 95\,mm baseline, $87^\circ\times58^\circ$ depth FOV, ideal range 0.6--6\,m and $<2\%$ depth error at 4\,m \cite{d455}.

The pipeline renders scenes in Blender, stores RGB, ideal $z$-depth, normals, instance and class labels and metadata, and then degrades the ideal depth with a cumulative, seeded ladder L0--L4 of stereo-inspired effects (Section~\ref{sec:sensor}). To show whether the camera model changes a downstream step, a classical localization baseline is run on every level (Section~\ref{sec:exp}). \emph{No real D455 data were used}: every parameter is a physically motivated assumption, so this is a controlled perturbation generator, not a calibrated simulator. The pipeline trains no model; domain randomization is complementary \cite{tobin2017,tremblay2018}.

\section{Pipeline}
\paragraph{Scenes.}
T-slot profiles (20/30/40\,mm sections, 0.10--0.40\,m long; silver anodised, black anodised or brushed aluminium) and L-brackets (30\,mm legs, 3\,mm thick, two 6\,mm holes per leg; zinc) are generated by Blender-independent geometry code. Each scene has 4--12 parts (bracket probability 0.35), dropped by rigid-body simulation into a $0.5\times0.4$\,m bin on a 2\,m table, under 1--3 randomized area lights. Material base colour, metallic, roughness and anisotropy are sampled from configured ranges. Units are SI.

\paragraph{Camera and rendering.}
A pinhole camera with $848\times480$ pixels, $87^\circ$ HFOV, 95\,mm baseline and 0.4--6.0\,m validity range ($f_x=f_y=\frac{W/2}{\tan(\mathrm{HFOV}/2)}$) is placed 0.45--0.9\,m above the bin, tilted 0--45$^\circ$ with random roll (eye-in-hand). Only baseline and HFOV are aligned to the datasheet; resolution and near cutoff are my choices. Geometry passes use one sample per pixel with no denoising and a near-zero filter, so labels, normals and depth do not mix at boundaries; colour uses 32 samples with denoising. Ray distance is converted to $z$-depth (0 = invalid). Every sample stores the format described above, including $K$, object poses and visible-pixel counts. All randomness derives from the config seed, and each sensor effect has its own random stream so removing one effect leaves the others' draws unchanged. The dev split is used for debugging and tuning; the test split (100 scenes, 793 instances: 522 profiles, 271 brackets) is held out.

\begin{figure}[tbp]
\centering
\includegraphics[width=0.78\textwidth]{outputs/montage.png}
\caption{Dev-split samples: RGB, colourized depth, instance labels, normals. Qualitative only.}
\label{fig:montage}
\end{figure}

\section{Camera model: included and omitted effects}
\label{sec:sensor}
The degradation operates on ideal $z$-depth, without Blender, through the rectified-stereo relation $d=f_xB/z$ ($d$ disparity in pixels, $B$ baseline). It does not simulate IR image formation or the D455's matching processor. Levels are cumulative (noise, quantization, dropout, then edges are applied in that order): L0 ideal; L1 + disparity quantization; L2 + disparity noise; L3 + dropout; L4 + occlusion/edge effects. Values outside 0.4--6.0\,m are invalid.

\textbf{Quantization (L1).} $d_q=\Delta\,\mathrm{round}(d/\Delta)$ with $\Delta=2^{-6}$\,px, giving depth steps that grow with range.
\textbf{Noise (L2).} $d_n=d+\epsilon$, $\epsilon\sim\mathcal N(0,\sigma_d^2)$, $\sigma_d=0.08\,s$\,px ($s=1$ nominal), applied before quantization. First-order, $\sigma_z\simeq z^2\sigma_d/(f_xB)$: 0.38\,mm at 0.45\,m, 1.53\,mm at 0.9\,m, 30\,mm at 4\,m (0.75\% of range, same order as but not the same metric as the datasheet's $<2\%$).
\textbf{Coherent dropout (L3).} With incidence angle $\theta=\arccos|\mathbf n^{\mathsf T}\mathbf r|$ and roughness $\rho$,
\[
p(\theta,\rho)=\operatorname{clip}_{[0,1]}\!\Big[p_{\rm base}+\tfrac{p_{\rm graze}}{1+e^{-(\theta-\theta_0)/w}}+p_{\rm spec}(1-\rho)e^{-\theta^2/2\sigma_{\rm spec}^2}\Big],
\]
with $p_{\rm base}=0.01$, $p_{\rm graze}=0.30$, $p_{\rm spec}=0.20$, $\theta_0=70^\circ$, $w=6^\circ$, $\sigma_{\rm spec}=12^\circ$. The grazing term makes oblique faces drop out; the specular term makes smooth metal seen head-on drop out. To avoid salt-and-pepper holes, Gaussian noise ($\sigma=4$\,px) is rank-transformed to a uniform field and pixels with field value $<p$ are invalidated, producing blobs.
\textbf{Edges (L4).} (i) Stereo occlusion: each pixel is projected to $u_r=\mathrm{round}(u-f_xB/z)$ into a right-image z-buffer and invalidated if it falls outside or lies more than 2\% behind the buffer. (ii) At depth steps $>0.05$\,m, pixels in a one-pixel band become flying pixels with probability 0.30 (uniform between the local 3$\times$3 min and max). (iii) Discontinuity pixels take the local 3$\times$3 minimum (foreground fattening).

\begin{table}[ht]
\centering
\footnotesize
\begin{tabular}{@{}p{3.9cm}p{1.9cm}p{10.5cm}@{}}
\toprule
Effect & Status & Reason \\
\midrule
Disparity quantization, $z^2$ noise & Included & Direct consequence of sub-pixel stereo; dominant, well-understood error trend. $\sigma_d$ is assumed. \\
Holes on shiny/oblique aluminium & Included & The most characteristic failure on anodised/brushed profiles and zinc brackets; modelled as a coherent incidence/roughness proxy. \\
Stereo occlusion, flying pixels, fattening & Included & Occlusion is geometry given the 95\,mm baseline; thin profile edges are where the others matter. Heuristic rates. \\
IR projector, pattern matching, specular/IR reflectance & Omitted & Needs an IR renderer and measured material response; the dropout proxy stands in for the net effect. \\
Exposure, gain, saturation, firmware filters & Omitted & Depend on settings and hardware measurements I do not have. \\
Calibrated bias, range-specific error, row/checkerboard hole patterns & Omitted & No measurements justify a pattern or values. \\
Motion/temporal effects, multipath & Omitted & Static scenes (global shutter); multipath needs full IR light transport. \\
\bottomrule
\end{tabular}
\caption{Included and omitted effects. All included parameters are modelling choices, not D455 measurements.}
\label{tab:effects}
\end{table}

\section{Experiment: does the camera model change localization?}
\label{sec:exp}
\paragraph{Design.}
The downstream step is per-instance 6-DoF localization in \emph{oracle-mask} mode: valid depth pixels inside the ground-truth mask are back-projected, cleaned by statistical outlier removal, initialized by PCA and refined by trimmed point-to-plane ICP against a seeded model cloud. Oracle masks isolate depth effects; this is not detection or grasping. Success: profiles lateral error $<5$\,mm and axis error $<5^\circ$; brackets centroid $<5$\,mm and symmetry-aware rotation $<10^\circ$. The baseline and thresholds were never changed. Arms: L0--L4, four leave-one-effect-out L4 variants, four single-effect arms, three edge components and noise scales $s\in\{0.5,1,2,4,8\}$ ($s=1$ reuses L2). Intervals are 95\% scene-bootstrap (2,000 replicates, seed 0, shared across arms so differences are paired). Arms and analyses were committed and tagged \texttt{v2-protocol} before the test run.

\paragraph{Main result.}
Success falls from 87.6\% (695/793) at L0 to 83.6\% (663/793) at L4. The paired L4--L0 change is $-4.04$ points overall (95\% CI $[-6.16,-2.10]$), $-1.72$ for profiles ($[-3.56,0.00]$) and $-8.49$ for brackets ($[-13.39,-3.91]$): the camera model changes the outcome, mainly for brackets.

\begin{table}[ht]
\centering
\footnotesize
\begin{tabular}{@{}lccccc@{}}
\toprule
Arm & All, \% (95\% CI) & Profile, \% & Bracket, \% (95\% CI) & Bracket pos. med/p90 & Bracket rot. med/p90 \\
\midrule
L0 & 87.6 (85.2--89.9) & 96.7 & 70.1 (64.1--75.9) & 3.20/8.74\,mm & 5.95/40.7$^\circ$ \\
L1 & 87.3 (84.7--89.6) & 96.7 & 69.0 (62.6--75.0) & 3.20/8.92\,mm & 6.88/46.2$^\circ$ \\
L2 & 86.0 (83.8--88.2) & 94.4 & 69.7 (63.9--75.2) & 3.07/9.10\,mm & 6.13/54.0$^\circ$ \\
L3 & 84.2 (81.9--86.6) & 94.6 & 64.2 (58.3--69.9) & 3.18/9.64\,mm & 7.60/61.6$^\circ$ \\
L4 & 83.6 (81.1--86.1) & 95.0 & 61.6 (55.4--67.7) & 3.29/11.03\,mm & 7.13/62.3$^\circ$ \\
\bottomrule
\end{tabular}
\caption{Localization success per level, with bracket error medians and 90th percentiles. Profile position median grows from 0.15\,mm (L0) to 0.24\,mm (L4). All 20 arms are in \texttt{outputs/v2/*.csv}.}
\label{tab:main}
\end{table}

\paragraph{Which effect matters.}
Degradation relative to L0 for single-effect arms (overall / bracket, points): quantization 0.38 / 1.11 (not detectable); noise 2.52 ($[0.74,4.39]$) / 2.21 ($[-1.29,5.79]$); dropout 1.77 ($[0.60,3.05]$) / 4.43 ($[1.38,7.53]$); bundled edges 1.13 ($[-0.26,2.59]$) / 3.32 ($[0.68,6.25]$). Within L4, removing dropout recovers 1.77 points overall and 4.43 for brackets; other leave-one-out intervals span zero. The L4 penalty is not detectably superadditive (residual $-1.77$, $[-4.32,0.79]$). Edge components alone (all / bracket success, \%): occlusion 85.9 / 65.7, flying pixels 88.4 / 70.1, fattening 88.5 / 72.0, so occlusion invalidation dominates the edge effect.

\paragraph{Noise sweep.}
Overall success is 86.8, 86.0, 82.7, 76.8 and 47.9\% at $s=0.5,1,2,4,8$; profile success falls monotonically (96.0 to 42.9\%), while bracket success is 69.0, 69.7, 74.2, 76.0 and 57.6\%. The bracket rise at $s=2$--4 is unexplained (likely noise helping ICP escape bad initializations, not tested), so no single noise-sensitivity claim is made.

\paragraph{Holes.}
The invalid fraction inside object masks rises from 9.3\% (L0--L2; due to range and background) to 14.2\% (L3) and 16.6\% (L4), while the connected invalid-region median size drops from 1,694\,px (L0) to 12\,px (L3) and 2\,px (L4): dropout and edges fragment large regions into many small holes (4,002 and 9,312 components versus 65).

\begin{figure}[tbp]
\centering
\begin{minipage}[t]{0.49\textwidth}\centering
\includegraphics[width=\linewidth]{outputs/v2/success_by_arm_v2.png}
\end{minipage}\hfill
\begin{minipage}[t]{0.49\textwidth}\centering
\includegraphics[width=\linewidth]{outputs/v2/noise_sweep_v2.png}
\end{minipage}
\caption{Success across all arms (left); success and p90 position error over the noise sweep (right).}
\label{fig:sweeps}
\end{figure}

\begin{figure}[tbp]
\centering
\includegraphics[width=0.78\textwidth]{outputs/v2/degradation_ladder_v2.png}
\caption{Test scene 00000, L0--L4 depth and absolute error relative to L0 (shared colour limits).}
\label{fig:ladder}
\end{figure}

\section{What did not work, and caveats}
\begin{itemize}
\setlength{\itemsep}{0pt}
\item \textbf{Weak baseline.} At L0, 98 of 793 instances already fail (81/271 brackets, 17/522 profiles): 65 ICP divergence/bad initialization, 28 orientation/symmetry ambiguity, 5 too few points. Success is 16\% below 200 valid points and above 97\% above 1,000. Bracket conclusions are therefore partly baseline-limited; I deliberately did not tune the baseline.
\item \textbf{Quantization is a null at this scale}, and in the full L4 context removing noise, quantization or edges has no detectable effect. This is conditional on the arm and class (noise alone is detectable, and the sweep is clear from $s=2$), not proof that these effects are irrelevant.
\item \textbf{Dropout parameter was set by eye.} The \texttt{v1-freeze} config had $p_{\rm graze}=0.60$; I lowered it to 0.30 after visually inspecting perturbed \emph{dev} maps because the holes looked implausibly large, so the freeze rule was broken for this one parameter. Test metrics were not seen; no other value was evaluated, so the dropout sensitivity is unknown. The decision date is not recorded.
\item \textbf{Working range.} 36.6\% of valid test pixels lie below the D455's 0.6\,m ideal minimum (scenes are 0.45--0.9\,m), so the model operates partly outside the datasheet's recommended range, where real error behaviour is not documented.
\item \textbf{Not run.} Spatially correlated noise (optional), detector-based masks and grasp success. The first script launch failed because Open3D was not on the import path; the rerun set the path and changed no data or parameter.
\item \textbf{No validation} against real D455 data: absolute numbers are not D455 performance.
\end{itemize}

\section{Conclusions}
The pipeline generates labelled aluminium-part RGB-D scenes with a seeded D455-inspired degradation ladder, replacing ideal depth. Quantization, range-growing noise, coherent incidence/roughness-dependent holes, stereo occlusion, flying pixels and fattening are modelled; IR/projector physics, exposure, calibrated bias, hole patterns, temporal effects and firmware filters are omitted because they need hardware measurements. The camera model does change the downstream outcome: oracle-mask success drops from 87.6\% to 83.6\% (paired $-4.04$ points, CI $[-6.16,-2.10]$), mostly for brackets ($-8.49$), with dropout the largest single bracket effect and noise the largest overall. Because every parameter is assumed, the experiment shows \emph{which} imperfections deserve calibration, not how the real camera performs.

\clearpage
\section*{Riskiest assumption for synthetic-to-real transfer, and a six-month plan}
\textbf{Riskiest assumption.} That the synthetic data cover the \emph{sensor failure modes} the real robot will meet, in particular depth on shiny, thin metal parts. A model trained on synthetic data sees only the artefacts someone chose to simulate, and on aluminium the dominant real artefacts (specular IR dropout, edge errors, projector interactions) depend on material, incidence and exposure in ways I could only guess at. My experiment makes this concrete: the dropout proxy was the largest bracket effect (4.43 points), yet its key parameter was set by eye, and halving it already changed how holes look. If the proxy's hole rates or sizes are wrong, a model can appear robust in simulation and fail in the cell, and a typical factory also adds unmodelled cell-specific effects (lighting, vibration, calibration drift). Evaluating only on synthetic data hides this gap.

\textbf{Months 1--2: measure the gap.} Fix camera firmware, stream mode, exposure, gain, emitter and calibration. Capture synchronized depth and left/right IR of profiles and brackets on a fixture with known poses across range, material, incidence and lighting, with a small labelled real test set that is never used for training or tuning.

\textbf{Months 3--4: quantify and calibrate the sensor model.} Compare invalid-pixel rate and connected-hole size/location distributions with geometric ground truth on held-out condition cells, before fitting anything. Go/no-go: median absolute error in invalid fraction $\le5$ points, and median and p90 hole areas within a factor of two of the real ones in at least 80\% of cells. If not met, do not claim transfer; replace the proxy with a fitted or learned (e.g.\ measured-IR) noise model.

\textbf{Months 5--6: test transfer.} Freeze the supported model, train or evaluate a perception model on synthetic data (with domain randomization over the calibrated parameters) and report real-set performance against a synthetic-only baseline. Ablate sensor effects as done here, to see which gap items drive real failures. Success is a real-set drop of at most a pre-stated margin; otherwise the finding and the failing condition cells are reported.

\section*{AI-use note}
I used AI assistants to implement the pipeline and experiment from prompts I wrote, to draft and tighten this report, and to review results and drafts. All numbers come from the generated outputs. \textbf{Disagreement/correction: [fill in one real case; do not invent one].}

\begin{thebibliography}{9}
\footnotesize
\bibitem{tobin2017} J. Tobin et al., ``Domain Randomization for Transferring Deep Neural Networks from Simulation to the Real World,'' IROS, 2017.
\bibitem{tremblay2018} J. Tremblay et al., ``Training Deep Networks with Synthetic Data: Bridging the Reality Gap by Domain Randomization,'' CVPR Workshops, 2018.
\bibitem{d455} Intel RealSense, \emph{D400 Series Datasheet} (D455: 95\,mm baseline, global shutter, $87^\circ\times58^\circ$ depth FOV, up to $1280\times720$, ideal range 0.6--6\,m, $<2\%$ error at 4\,m), 2023. \url{https://www.realsenseai.com/wp-content/uploads/2023/07/Intel-RealSense-D400-Series-Datasheet-July-2023.pdf}
\end{thebibliography}
\end{document}
